% Source-derived chapter 03: Non-unique specializations of ramified Galois sections
% Original source: birational-section-conjecture-strong-birationality
\chapter{Non-unique specializations of ramified Galois sections}
\label{birational-section-conjecture-strong-birationality-chapter-03}
\source{birational-section-conjecture-strong-birationality}

\begin{remark}[Source section]
\label{birational-section-conjecture-strong-birationality-chapter-03-source}
\source{birational-section-conjecture-strong-birationality}
\source{birational-section-conjecture-strong-birationality}
This chapter reproduces the corresponding section of the retrieved
LaTeX source, including its definitions, statements, examples, and
proofs. It has not yet been translated into Lean.
\notready
\end{remark}

% The source section title was: Non-unique specializations of ramified Galois sections
\label{ramspe}
\source{birational-section-conjecture-strong-birationality}

Classically, the specialization of a Galois section is defined if the original Galois section is \emph{unramified} \cite[\S 8.2]{sti13}, i.e. when the \emph{ramification homomorphism} is trivial. A notion of specialization exists \emph{always}, as long as we don't require that the specialization is unique. If the section is unramified, this "generalized specialization" is unique and coincides with the classical specialization. While a definition of specialization for ramified sections exists in the literature \cite[Corollary 8.9]{sti13}, little else in known beyond the definition itself. 



%~ \subsection{Specializing loops}

%~ Let $R$ be a DVR with fraction field $K$ and residue field $k$ and $X\to \spec R$ a proper morphism of schemes. A generic section $\spec K\to X_{K}$ extends uniquely to a section $\spec R\to X$ thanks to the valuative criterion of properness, this gives a specialization $\spec k\to X_{k}$. Suppose now that $X\to\spec R$ is a family of curves, then $\Pi_{X/R}\to\spec R$ is a projective limit of proper morphisms by \autoref{limit}: given a generic Galois section $\spec K\to\Pi_{X_{K}/K}$, we would like to find a specialization. 



%~ The classical valuative criterion for proper morphisms of stacks \cite[Théorème 7.3]{lmb00} does not help us, since it requires passing to an extension of $R$ and thus enlarging the residue field. 

Let $R$ be a DVR with fraction field $K$ and residue field $k$, the \emph{$n$-th root stack} $\sqrt[n]{\spec R}$ of $\spec R$ is defined for every $n$ \cite[Appendix B]{agv08}. If $\pi\in R$ is a uniformizing parameter, then $\sqrt[n]{\spec R}$ is isomorphic to the quotient stack $[\spec R(\sqrt[n]{\pi})/\mu_{n}]$, but the definition of root stack does not depend on the choice of $\pi$. The morphism $\sqrt[n]{\spec R}\to \spec R$ is generically an isomorphism, while the closed fiber with the reduced structure is non-canonically isomorphic to the classifying stack $B\mu_{n}$.

The \emph{infinite root stack} $\sqrt[\infty]{\spec R}$ is the projective limit $\projlim_{n}\sqrt[n]{\spec R}$, see \cite{tv18} for details. Let $c\in\spec R$ be the closed point and $S^{1}_{c}$ the reduced fiber of $\sqrt[\infty]{\spec R}$ over $c$, it is non-canonically isomorphic to the classifying stack $\cB_{k}\hz(1)$ of $\hz(1)$ over $k$. The notation $S^{1}_{c}$ is meant to be reminiscent of the topological space $S^{1}$, which is the classifying space $\mathcal{B}\Z$ of $\Z$. There is a non-canonical isomorphism between the isomorphism classes of $S^{1}_{c}(k)$ and $\H^{1}(k,\hz(1))=\projlim_{n}k^*/k^{*n}=\hat{k^*}$.
 %~ Because of this, the correct topological analogue of $\cB_{k(c)}\hz(1)$ is $\C^{*}$ rather than $S^{1}$, but using the notation $\C^{*}_{c}$ would be extremely confusing.

Recall that a morphism $Y\to X$ of fibered categories over $k$ is \emph{constant} if there exists a factorization $Y\to\spec k\to X$.

\begin{proposition}
\source{birational-section-conjecture-strong-birationality}
\notready\label{prorootval}
\source{birational-section-conjecture-strong-birationality}
	Let $X\to C$ be a family of curves as in \autoref{families}, $c\in C$ a closed point with local ring $R$ and $z:\spec K\to\Pi_{X/C}$ a generic Galois section, where $K$ is the fraction field of $C$. Then $z$ extends to a $2$-commutative diagram
	\[\begin{tikzcd}
		\spec K\rar\ar[dr]		&	\sqrt[\infty]{\spec R}\rar[dotted]\dar	&	\Pi_{X/C}\dar	\\
								&	\spec R\rar												&	C.
	\end{tikzcd}\]
	The extension $\sqrt[\infty]{\spec R}\to\Pi_{X/C}$ is unique up to a unique isomorphism. 
	
	We call the induced morphism $S^{1}_{c}\to\Pi_{X_{c}/k(c)}$ the specializing loop $\gamma_{z}(c)$ of $z$ at $c$. An extension $\spec R\to\Pi_{X/C}$ exists if and only if the specializing loop is constant.
\end{proposition}

\begin{proof}
	Using \autoref{limit}, this is a direct consequence of \cite[Theorem 3.1]{giulio-angelo-valuative}.
\end{proof}

%~ The name ``specializing loop'' is meant to be reminiscent of the fact that the circle $S^{1}$ is the classifying space $\mathcal{B}\Z$ of $\Z$ and a loop is a map from $S^{1}=\mathcal{B} \Z$ to some other space, just as the specializing loop is a map from $S^{1}_{c}\simeq\mathcal{B}_{k(c)}\hz(1)$ to the fundamental gerbe.



%~ \begin{theorem}\label{prorootval}
%~ Let $f:\mc{X}=\projlim_{i}\mc{X}_{i}\to \spec R$ be a projective limit of proper morphisms of Deligne-Mumford stacks, where $R$ is a DVR of pure characteristic $0$. Let $c\in\spec R$ be the closed point, $K$ the fraction field and $k$ the residue field. Suppose that we have a generic section $\spec K\to \mc{X}$. Then there exists an extension $\sqrt[\infty]{\spec R} \to \mc{X}$, unique up to a unique isomorphism, making the diagram
   %~ \[\begin{tikzcd}
   %~ \spec K\rar\ar[d, hook]  						&   \mc{X}\dar   \\
   %~ \sqrt[\infty]{\spec R}\rar\ar[ru,"\exists !"]   	&   \spec R
   %~ \end{tikzcd}\]
%~ $2$-commutative. Moreover, an extension $\spec R\to\mc{X}$ exists if and only if the induced morphism $H_{c}\to\mc{X}_{k}$ is constant.
%~ \end{theorem}

%~ The proof of a more general form of \autoref{prorootval} in arbitrary characteristic will appear in a joint article with A. Vistoli \cite{giulio-angelo-lang-nishimura}. In the meantime, a proof in characteristic $0$ can be found in a previous version of this article \cite[version 1, Appendix A]{tbir1}.

%~ Let $C,X$ be noetherian, normal, connected schemes and $X\to C$ a geometric fibration with connected fibers. Let $z:\spec k(C)\to\Pi_{X/C}$ be a generic Galois section, where $k(C)$ is the fraction field of $C$. For every codimension $1$ point $c\in C$ the valuative criterion \autoref{prorootval} applied to $X_{\O_{C,c}}\to\spec \O_{C,c}$ induces a morphism $h_{z}(c):H_{c}\to\Pi_{X_{k(c)}/k(c)}$: we call this \emph{the specializing loop} of $z$ at $c$. If $C$ is the spectrum of a DVR, we may just write $h_{z}$.

\begin{definition}
\source{birational-section-conjecture-strong-birationality}
\notready\label{loopdef}
\source{birational-section-conjecture-strong-birationality}
	%~ With notation as in \autoref{prorootval}, the composition $H_{c}\to\sqrt[\infty]{\spec R}\to\Pi_{X/C}$ factorizes as a morphism
	%~ \[H_{c}\to\Pi_{X_{c}/k(c)}\]
	%~ to the fundamental gerbe of the fiber, we call this morphism the \emph{specializing loop} $h_{z}(c)$ of $z$ at $c$. 
	A \emph{specialization} of $z$ at $c$ is any Galois section of $X_{c}/k(c)$ in the essential image of the specializing loop $S^{1}_{c}(k(c))\to\Pi_{X_{c}/k(c)}(k(c))$ of $z$ at $c$ evaluated on $k(c)$. A specialization always exists, but in general it might be not unique. If the specializing loop is constant, then the specialization is unique (the converse is false in general).
\end{definition}

If we fix any section $s\in S^{1}_{c}(k(c))$, the specializing loop $\gamma_{z}:S^{1}_{c}\to\Pi_{X_{c}/k_{c}}$ is constant if and only if the corresponding homomorphism of group schemes $\hz(1)=\uaut(s)\to\uaut(\gamma_{z}(s))$ is trivial. Suppose for simplicity that $k(c)=k$. By the characterization given above, the specializing loop of $z$ at $c$ is constant if and only if the specializing loop of $z_{\bar{k}(C)}$ is constant at $c_{\bar{k}}$: having constant specializing loop in the closed fiber $X_{c}$ only depends on the base change to $\bar{k}$ of the family, nonetheless it encapsulates arithmetic information over $k$, i.e. uniqueness of the specialization.

%~ \begin{lemma}\label{uniqueloop1}
	%~ Let $X\to \spec R$ a geometric fibration, $z\in\Pi_{X_{K}/K}(K)$ a generic section, $k'/k$ an extension of the residue field. Then the specializing loop $\gamma_{z}$ is constant if and only if the base change $\gamma_{z,k'}$ is constant.
	%~ \begin{proof}
		%~ Choose a section of $H_{c}(k)$ so that we may identify $H_{c}=B\hz(1)$, $\Pi_{X_{k}/k}=BG$ for some group scheme $G$ over $k$. The specializing loop is constant if and only if the associated homomorphism $\hz(1)\to G$ is trivial, and this condition is clearly invariant under base change.
	%~ \end{proof}
%~ \end{lemma}

\begin{remark}
\source{birational-section-conjecture-strong-birationality}
\notready
	A specialization of a ramified section can be constructed in the language of étale fundamental groups as follows (see also \cite[Corollary 89]{sti13}). Let $R$ be a DVR with fraction and residue fields $K$, $k$ of characteristic $0$, denote by $R^{h}$ a henselianization and $K^{h}$ the fraction field. We may identify $\gal(\bar{k}/k)=\pi_{1}(\spec R^{h})$, and if $X\to \spec R$ is a family of curves then $\pi_{1}(X_{k})\simeq\pi_{1}(X_{R^{h}})$. Since $R^{h}$ is henselian, the Galois group $\gal(\bar{K}/K^{h})$ coincides with the decomposition subgroup and hence it is an extension of $\gal(\bar{k}/k)$ by the inertia subgroup.
	
	Let $\pi\in R^{h}$ be a uniformizing parameter, choose a $n$-th root of $\pi$ for every $n$ in a compatible way and let $K^{h}(\sqrt[\infty]{\pi})$ be the corresponding extension of $K^{h}$, we have that $\gal(\bar{K}/K^{h}(\sqrt[\infty]{\pi}))\simeq\gal(\bar{k}/k)$ defines a section of $\gal(\bar{K}/K^{h})\to\pi_{1}(\spec R^{h})=\gal(\bar{k}/k)$.
	 
	
	If $\gal(\bar{K}/K)\to\pi_{1}(X_{K})$ is a generic Galois section, it induces a Galois section $\gal(\bar{K}/K^{h})\to\pi_{1}(X_{K^{h}})$ and a specialization is defined by the composition
	\[\gal(\bar{k}/k)\to \gal(\bar{K}/K^{h})\to \pi_{1}(X_{K^{h}})\to\pi_{1}(X_{R^{h}})\simeq\pi_{1}(X_{k}).\]
	Let $I\subset \gal(\bar{K}/K^{h})$ be the inertia subgroup, by composition we have a map $I\to\pi_{1}(X_{\bar{k}})\subset\pi_{1}(X_{k})$ which is called \emph{ramification map}: unramified Galois sections are those for which this map is trivial. 
	
	If we fix a section of $D\to\gal(\bar{k}/k)$, the corresponding specialization induces an action of $\gal(\bar{k}/k)$ on $\pi_{1}(X_{\bar{k}})$, and the ramification map is Galois-equivariant with respect to the natural Galois action on $I\simeq\hz(1)$. The set of specializations is the image in nonabelian Galois cohomology of the ramification map. 
	
	%~ Furthermore, for our arguments it would be necessary to define the specializing loop as a Galois equivariant homomorphism from the inertia group to the étale fundamental group of the fiber: again, this depends on some choices, while the specializing loop as a morphism of gerbes is canonical.
\end{remark}

\begin{example}
\source{birational-section-conjecture-strong-birationality}
\notready\label{gmloop}
\source{birational-section-conjecture-strong-birationality}
	Let $R$ be a DVR as in \autoref{families} with fraction field $K$ and valuation $v:K^{*}\to\Z$, this extends naturally to an homomorphism $\hat{K^{*}}\to\hz$ which we still call $v$. The morphism $\G_{m}=\A^{1}_{R}\setminus\{0\}\to \spec R$ is a family of curves.
	
	The section $1\in\G_{m}(R)$ gives an identification $\Pi_{\G_{m}/R}=\cB_{R}\hz(1)$, in particular $\Pi_{\G_{m}/R}(R)=\hat{R^*}$ and $\Pi_{\G_{m,K}/K}(K)=\hat{K^*}$. If $z\in\hat{K^{*}}$ is a generic section, by \autoref{prorootval} the specializing loop $\gamma_{z}$ is constant if and only if $z$ is in the image of $\hat{R^{*}}\to\hat{K^{*}}$, or equivalently if and only if $v(z)=0\in\hz$.
\end{example}

It can be proved that unramified Galois sections in the sense of \cite[Chapter 8]{sti13} are those for which the specializing loop is constant, though we don't need this fact. If the specializing loop is constant, the specialization is clearly unique; this is coherent with the fact that unramified Galois sections have a canonical specialization.

The converse is false in general: if the residue field is algebraically closed, the specialization is unique since there is only one $\hz(1)$-torsor up to equivalence, but the Galois section might be ramified and the specializing loop not constant. For instance, we may choose $R=\C[x]_{(x)}$ in \autoref{gmloop}, the generic Galois section associated with $x\in\hat{\C(x)^{*}}\setminus\hat{\C[x]_{(x)}^{*}}$ has non-constant specializing loop but unique specialization. 
%~ The bottom line is that ramification is a geometric concept, while uniqueness of specialization is an arithmetic one.

Still, in arithmetic situations, it is often the case that uniqueness of specialization forces a constant specializing loop. 

\begin{lemma}
\source{birational-section-conjecture-strong-birationality}
\notready\label{uniqueloop}
\source{birational-section-conjecture-strong-birationality}
	Let $T$ be a torus over a field $k$ with a surjective valuation $\nu:k^{*}\to\Z$. Suppose that we have a morphism of gerbes $\gamma:\cB_{k}\hz(1)\to \Pi_{T/k}$. If the image of $\cB_{k}\hz(1)(k)$ in $\Pi_{T/k}(k_{\nu})$ has a finite number of isomorphism classes, then $\gamma$ is constant.
 %, or equivalently the induced homomorphism $f:\hz(1)\to\upi_{1}(T)$ is trivial.
	\begin{proof}
		%~ We may identify $\Pi_{T/k}$ with $\cB_{k} G$ where $G=\projlim_{n}T[n]$ is the Tate module of $T$, then $f$ is associated with an homomorphism $g:\hz(1)\to G$. We have that $f$ is constant if and only if $g$ is trivial, or equivalently if and only if
    Using the image of the preferred $k$-rational section of $\cB_{k}\hz(1)$ corresponding to the trivial torsor, we may identify $\Pi_{T/k}$ with $\mc{B}_{k}\upi_{1}(T)$, and we reduce to prove the following: if $f:\hz(1)\to\pi_{1}(T_{\bar{k}})$ is a Galois equivariant homomorphism such that the composition $\H^{1}(k,\hz(1))\xar{f} \H^{1}(k,\pi_{1}(T_{\bar{k}})) \xar{res} \H^{1}(k_{\nu},\pi_{1}(T_{\bar{k}}))$ has finite image, then $f$ is trivial. 

    Choose $k'/k$ a finite extension such that $T_{k'}\simeq\G_{m}^{n}$, we get an identification $\pi_{1}(T_{\bar{k}})\simeq \hz^{n}(1)$ of Galois modules over $k'$. Let $\nu'$ be an extension of $\nu$ to $k'$, denote by $r$ the ramification index of $\nu'/\nu$. The valuation $\nu$ induces an homomorphism $\H^{1}(k_{\nu},\hz(1))=\projlim_{i}k_{\nu}^{*}/k_{\nu}^{*i}\to \hz$, and similarly for $k'_{\nu'}$; with an abuse of notation, we denote these homomorphisms by $\nu$, $\nu'$. 
    
    %The hypothesis implies that the composition 
    %\[\H^{1}(k,\hz(1)) \xar{f}  \H^{1}(k,\pi_{1}(T_{\bar{k}})) \xar{res} \H^{1}(k_{\nu},\pi_{1}(T_{\bar{k}})) \xar{res} \H^{1}(k'_{\nu'},\pi_{1}(T_{\bar{k}}))=\H^{1}(k'_{\nu'},\hz(1)^{n})\] has finite image.

    We have a commutative diagram
    \[\begin{tikzcd}
        \H^{1}(k,\hz(1))\rar["res"] \ar[dd,"res\circ \H^{1}(f)"]     &   \H^{1}(k_{\nu},\hz(1))\dar["res"]\rar["\nu"]    &   \hz\dar["r"]    \\
                                                &   \H^{1}(k'_{\nu'},\hz(1))\dar["\H^{1}(f)"]\rar["\nu'"]            &   \hz\dar["f"]         \\
        \H^{1}(k_{\nu},\pi_{1}(T_{\bar{k}}))\rar["res"] & \H^{1}(k'_{\nu'},\pi_{1}(T_{\bar{k}})\simeq\hz(1)^{n})\rar["\nu'"] & \pi_{1}(T_{\bar{k}})\simeq \hz^{n}
    \end{tikzcd}\]
    where the last identification $\pi_{1}(T_{\bar{k}})\simeq \hz^{n}$ ignores the Galois action. The hypothesis implies that the composition 
    \[f\circ r \circ \nu \circ res = \nu' \circ res \circ \H^{1}(f) : \H^{1}(k,\hz(1)) \to \hz^{n}\]
    has finite image, hence it is $0$ since $\hz$ has no torsion.  Notice that, since $\nu:k^{*}\to\Z$ is surjective, then the upper horizontal arrow $\nu \circ res: \H^{1}(k,\hz(1))=\projlim_{i}k^{*}/k^{*i} \to \hz$ is surjective as well, and hence $f\circ r =rf=0$. Since $r\neq 0$, we get that $f=0$, as desired.
    
%  It is enough to show that $f_{\ell}:\Z_{\ell}(1)\to \on{T}_{\ell}T=\projlim_{n}T[\ell^{n}]$ is trivial for every prime $\ell$. There exists a finite extension $h/k$ such that $T_{h}=\G_{m}^{n}$, let $f_{i}:\Z_{\ell}(1)\to\Z_{\ell}(1)$ be the $i$-th coordinate of $f_{\ell,h}$, it is enough to show that $f_{i}$ is trivial for every $i$.
		
%		Let $\mu$ be an extension of $\nu$ to $h$, $r$ the ramification index of $\mu/\nu$. Denote by $I_{\ell},J_{\ell}$ the projective limits $\projlim_{n}k^{*}/k^{*\ell^{n}}$, $\projlim_{n}h_{\mu}^{*}/h_{\mu}^{*\ell^{n}}$, we have $\H^{1}(k,\Z_{\ell})=I_{\ell}$, $\H^{1}(k_{\nu},\Z_{\ell})=J_{\ell}$ and the valuations $\nu,\mu$ induce a natural commutative diagram
%		\[\begin{tikzcd}
%			I_{\ell}\rar["\nu"]\dar	&	\Z_{\ell}\dar["r"]	\\
%			J_{\ell}\rar["\mu"]		&	\Z_{\ell}
%		\end{tikzcd}\]
		
%		If $\phi:J_{\ell}\to J_{\ell}$ is the homomorphism induced by $f_{i}$, the image of the composition
%		\[I_{\ell}\to J_{\ell}\xar{\phi}J_{\ell}\xar{\mu}\Z_{\ell}\]
%		is $rf_{i}(1)\Z_{\ell}$. The hypothesis implies that $I_{\ell}\to J_{\ell}\xar{\phi}J_{\ell}$ has finite image, hence $f_{i}(1)=0$ and $f_{i}$ is trivial.
	\end{proof}
\end{lemma}

\begin{lemma}
\source{birational-section-conjecture-strong-birationality}
\notready\label{uniqueloop1}
\source{birational-section-conjecture-strong-birationality}
	Let $k$ be a number field with a finite place $\nu$, $X$ a curve over $k$, $\gamma:\cB_{k}\hz(1)\to\Pi_{X/k}$ a morphism of gerbes. If the image of $\cB_{k}\hz(1)(k)$ in $\Pi_{X/k}(k_{\nu})$ has a finite number of isomorphism classes, then $\gamma$ is constant.
	\begin{proof}
		If by contradiction $\gamma$ is not constant, i.e. the associated homomorphism $\hz\to\pi_{1}(X_{\bar{k}})$ is not trivial, up to replacing $X$ with a finite étale cover we may assume that the composition $\cB_{k}\hz(1)\to\Pi_{X/k}\to\Pi_{X/k}^{\rm ab}$ is not constant, where $\Pi_{X/k}^{\rm ab}$ is the abelianized fundamental gerbe, equivalently we may assume that the induced homomorphism $\hz\to\pi_{1}(X_{\bar{k}})^{\rm ab}$ is not trivial. We may do so because the map of Galois sections associated with a finite étale cover has finite fibers, hence the hypothesis remains true after passing to the covering. 
		
		Let $J$ be the semi-abelian Jacobian of $X$, it is an extension of an abelian variety $A$ by a torus $T$ and we have a canonical identification $\Pi_{X/k}^{\rm ab}=\Pi_{J/k}$. We may use the images of preferred section of $\cB_{k}\hz(1)$ to give identifications $\Pi_{J/k}=\cB_{k}\on{T}J$, $\Pi_{A/k}=\cB_{k}\on{T}A$ where $\on{T}J=\projlim_{n}J[n],\on{T}A=\projlim_{n}A[n]$ are the Tate modules; the morphisms of gerbes then correspond to homomorphisms $\hz(1)\to\on{T}J$, $\on{T}J\to\on{T}A$.
		

		
		 %~ $\cB_{k}\hz(1)\to\Pi_{T/k}\to\Pi_{J/k}$ since $\Pi_{A/k}=\cB_{k}\on{T}A$.
		By weight reasons, the composition $\hz(1)\to\on{T}J\to\on{T}A$ is trivial, hence we have a factorization $\hz(1)\to\on{T}T\to\on{T}J$. We have a short exact sequence
		\[\projlim_{n}A[n](k_{\nu})\to\H^{1}(k_{\nu},\on{T}T)\to\H^{1}(k_{\nu},\on{T}J).\]
		By Mattuck's theorem \cite{mat55} $A(k_{\nu})$ has finite torsion, thus $\H^{1}(k_{\nu},\on{T}T)\to\H^{1}(k_{\nu},\on{T}J)$ is injective. The hypothesis implies that the image of $\H^{1}(k,\hz(1))$ in $\H^{1}(k_{\nu},\on{T}J)$ is finite, hence the image in $\H^{1}(k_{\nu},\on{T}T)$ is finite too. By \autoref{uniqueloop}, $\hz(1)\to\on{T}T$ is trivial, which is contradiction with the fact that $\cB_{k}\hz(1)\to\Pi_{X/k}^{\rm ab}$ is not constant.
		
		%~ $B\hz(1)(k)/{\sim}$ in $\Pi_{J/k}(k_{\nu})/{\sim}$ is finite. We have a short exact sequence
		%~ \[\projlim_{n}A[n](k_{\nu})\to\H^{1}(k_{\nu},\on{T}T)\to\H^{1}(k_{\nu},\on{T}J).\]
		%~ By Mattuck's theorem \cite{mat55} $A(k_{\nu})$ has finite torsion, thus $\Pi_{T/k}(k_{\nu})\to\Pi_{J/k}(k_{\nu})$ is injective on isomorphism classes and thus the image of $B\hz(1)(k)/{\sim}$ in $\Pi_{T/k}(k_{\nu})/{\sim}$ is finite, too. This implies that $B\hz(1)\to\Pi_{T/k}$ is constant by \autoref{uniqueloop}, hence we have a contradiction.
	\end{proof}
\end{lemma}
